% GENERATED FILE. Edit canonical lessons, then run npm run generate:books.
% canonical-source-hash: fnv1a64:156db54db0395ae9
% canonical-lessons: MW-C196-dak, MW-C196-liphapho, MW-C196-parsal, MW-C196-thikano, MW-C196-mokalno

\chapter{At the Post}
\label{ch:mw-at-the-post}
\hlchaptermodality{196}

\begin{chapteropening}
\textbf{By the end of this chapter:} \emph{I can send a letter or a parcel.}

Complete the last lesson of chapter 196: I can send a letter or a parcel.
\end{chapteropening}

\section[\texorpdfstring{\mw{डाक} (ḍāk)}{डाक (ḍāk)}]{\mw{डाक} (ḍāk) — the post; mail}

\label{lesson:MW-C196-dak}



\begin{quote}
Before the new one: say the Marwadi for remaining, then the Marwadi for to be left over.
\end{quote}

\subsection*{You'll want to know: \mw{डाक}}
\textbf{\mw{डाक}} — \emph{ḍāk} — \textquotedblleft{}the post; mail\textquotedblright{}.

\begin{quote}
\mw{आज} \mw{डाक} \mw{कोनी} \mw{आई।}
\end{quote}

\begin{quote}
\emph{āj ḍāk konī āī.}
\end{quote}

\begin{quote}
The post didn't come today.
\end{quote}

A feminine noun.

\begin{cousinweb}
\textbf{\mw{डाक}} (\emph{ḍāk}) is inherited from a reconstructed \emph{*ḍakka}, \textquotedblleft{}relay, post\textquotedblright{}. CDIAL gives no Marwari form in this entry, but lists Hindi and Gujarati \emph{ḍāk}, \textquotedblleft{}relay, post, posting stage\textquotedblright{} (CDIAL 5519).
\end{cousinweb}

\begin{grammarlens}[title={What you've built}]
One more word for reading notices and messages.
\end{grammarlens}

\subsection*{Your turn}
\begin{itemize}
\raggedright
  \item \emph{Say it:} \emph{ḍāk}
  \item \emph{Say it:} \emph{ḍāk}, once more
  \item \emph{Say it:} \emph{āj ḍāk konī āī.} — The post didn't come today.
  \item \emph{Recall:} say the Marwadi for remaining, then the Marwadi for to be left over, then say \emph{ḍāk} again
\end{itemize}

\begin{tcolorbox}[breakable,colback=teal!4,colframe=teal!35!black,title={Before you move on}]
What does \mw{डाक} mean? (\textquotedblleft{}The post; mail\textquotedblright{}.) Say it once more.
\end{tcolorbox}
\section[\texorpdfstring{\mw{लिफाफो} (liphāpho)}{लिफाफो (liphāpho)}]{\mw{लिफाफो} (liphāpho) — an envelope}

\label{lesson:MW-C196-liphapho}



\begin{quote}
Before the new one: say the Marwadi for to be left over, then the Marwadi for the post.
\end{quote}

\subsection*{You'll want to know: \mw{लिफाफो}}
\textbf{\mw{लिफाफो}} — \emph{liphāpho} — \textquotedblleft{}an envelope\textquotedblright{}.

\begin{quote}
\mw{लिफाफो} \mw{खाली} \mw{है।}
\end{quote}

\begin{quote}
\emph{liphāpho khālī hai.}
\end{quote}

\begin{quote}
The envelope is empty.
\end{quote}

A masculine noun.

\begin{cousinweb}
\textbf{\mw{लिफाफो}} (\emph{liphāpho}) is usually traced, through Persian, to Arabic \emph{lifāfa}, \textquotedblleft{}a wrapper\textquotedblright{}.
\end{cousinweb}

\begin{grammarlens}[title={What you've built}]
One more word for reading notices and messages.
\end{grammarlens}

\subsection*{Your turn}
\begin{itemize}
\raggedright
  \item \emph{Say it:} \emph{liphāpho}
  \item \emph{Say it:} \emph{liphāpho}, once more
  \item \emph{Say it:} \emph{liphāpho khālī hai.} — The envelope is empty.
  \item \emph{Recall:} say the Marwadi for to be left over, then the Marwadi for the post, then say \emph{liphāpho} again
\end{itemize}

\begin{tcolorbox}[breakable,colback=teal!4,colframe=teal!35!black,title={Before you move on}]
What does \mw{लिफाफो} mean? (\textquotedblleft{}An envelope\textquotedblright{}.) Say it once more.
\end{tcolorbox}
\section[\texorpdfstring{\mw{पारसल} (pārsal)}{पारसल (pārsal)}]{\mw{पारसल} (pārsal) — a parcel}

\label{lesson:MW-C196-parsal}



\begin{quote}
Before the new one: say the Marwadi for the post, then the Marwadi for an envelope.
\end{quote}

\subsection*{You'll want to know: \mw{पारसल}}
\textbf{\mw{पारसल}} — \emph{pārsal} — \textquotedblleft{}a parcel\textquotedblright{}.

\begin{quote}
\mw{पारसल} \mw{घणो} \mw{भारी} \mw{है।}
\end{quote}

\begin{quote}
\emph{pārsal ghaṇo bhārī hai.}
\end{quote}

\begin{quote}
The parcel is very heavy.
\end{quote}

A masculine noun.

\begin{cousinweb}
\textbf{\mw{पारसल}} (\emph{pārsal}) is borrowed from English \emph{parcel}.
\end{cousinweb}

\begin{grammarlens}[title={What you've built}]
One more word for reading notices and messages.
\end{grammarlens}

\subsection*{Your turn}
\begin{itemize}
\raggedright
  \item \emph{Say it:} \emph{pārsal}
  \item \emph{Say it:} \emph{pārsal}, once more
  \item \emph{Say it:} \emph{pārsal ghaṇo bhārī hai.} — The parcel is very heavy.
  \item \emph{Recall:} say the Marwadi for the post, then the Marwadi for an envelope, then say \emph{pārsal} again
\end{itemize}

\begin{tcolorbox}[breakable,colback=teal!4,colframe=teal!35!black,title={Before you move on}]
What does \mw{पारसल} mean? (\textquotedblleft{}A parcel\textquotedblright{}.) Say it once more.
\end{tcolorbox}
\section[\texorpdfstring{\mw{ठिकाणो} (ṭhikāṇo)}{ठिकाणो (ṭhikāṇo)}]{\mw{ठिकाणो} (ṭhikāṇo) — an address; whereabouts}

\label{lesson:MW-C196-thikano}



\begin{quote}
Before the new one: say the Marwadi for an envelope, then the Marwadi for a parcel.
\end{quote}

\subsection*{You'll want to know: \mw{ठिकाणो}}
\textbf{\mw{ठिकाणो}} — \emph{ṭhikāṇo} — \textquotedblleft{}an address; whereabouts\textquotedblright{}.

\begin{quote}
\mw{म्हारो} \mw{ठिकाणो} \mw{चिट्ठी} \mw{माथै} \mw{है।}
\end{quote}

\begin{quote}
\emph{mhāro ṭhikāṇo chiṭṭhī māthai hai.}
\end{quote}

\begin{quote}
My address is on the letter.
\end{quote}

A masculine noun.

\begin{cousinweb}
\textbf{\mw{ठिकाणो}} (\emph{ṭhikāṇo}) is inherited from a reconstructed \emph{*ṭhīkka}, \textquotedblleft{}firm, right\textquotedblright{}; CDIAL lists Old Marwari \emph{ṭhīka}, \textquotedblleft{}trace, whereabouts\textquotedblright{}, and Hindi \emph{ṭhikānā} under it (CDIAL 5503).
\end{cousinweb}

\begin{grammarlens}[title={What you've built}]
One more word for reading notices and messages.
\end{grammarlens}

\subsection*{Your turn}
\begin{itemize}
\raggedright
  \item \emph{Say it:} \emph{ṭhikāṇo}
  \item \emph{Say it:} \emph{ṭhikāṇo}, once more
  \item \emph{Say it:} \emph{mhāro ṭhikāṇo chiṭṭhī māthai hai.} — My address is on the letter.
  \item \emph{Recall:} say the Marwadi for an envelope, then the Marwadi for a parcel, then say \emph{ṭhikāṇo} again
\end{itemize}

\begin{tcolorbox}[breakable,colback=teal!4,colframe=teal!35!black,title={Before you move on}]
What does \mw{ठिकाणो} mean? (\textquotedblleft{}An address; whereabouts\textquotedblright{}.) Say it once more.
\end{tcolorbox}
\section[\texorpdfstring{\mw{मोकलणो} (mokalṇo)}{मोकलणो (mokalṇo)}]{\mw{मोकलणो} (mokalṇo) — to send off}

\label{lesson:MW-C196-mokalno}



\begin{quote}
Before the new one: say the Marwadi for a parcel, then the Marwadi for an address.
\end{quote}

\subsection*{You'll want to know: \mw{मोकलणो}}
\textbf{\mw{मोकलणो}} — \emph{mokalṇo} — \textquotedblleft{}to send off\textquotedblright{}.

\begin{quote}
\mw{हूं} \mw{काल} \mw{पारसल} \mw{मोकलसूं।}
\end{quote}

\begin{quote}
\emph{hū̃ kāl pārsal mokalsū̃.}
\end{quote}

\begin{quote}
I'll send the parcel off tomorrow.
\end{quote}

Past (I …): \mw{म्हैं} \mw{चिट्ठी} \mw{मोकली}, \emph{mhaĩ chiṭṭhī moklī}.

\begin{cousinweb}
\textbf{\mw{मोकलणो}} (\emph{mokalṇo}) is inherited from a reconstructed \emph{*mukna}, \textquotedblleft{}loosed\textquotedblright{}; CDIAL lists the Old Marwari verb \emph{mokaï}, \textquotedblleft{}sends\textquotedblright{}, and Sindhi \emph{mokilaṇu}, \textquotedblleft{}to send\textquotedblright{}, with the \emph{l} this spelling keeps (CDIAL 10157): to send a thing is to let it go.
\end{cousinweb}

\begin{grammarlens}[title={What you've built}]
One more word for reading notices and messages.
\end{grammarlens}

\subsection*{Your turn}
\begin{itemize}
\raggedright
  \item \emph{Say it:} \emph{mokalṇo}
  \item \emph{Say it:} \emph{mokalṇo}, once more
  \item \emph{Say it:} \emph{hū̃ kāl pārsal mokalsū̃.} — I'll send the parcel off tomorrow.
  \item \emph{Recall:} say the Marwadi for a parcel, then the Marwadi for an address, then say \emph{mokalṇo} again
\end{itemize}

\begin{tcolorbox}[breakable,colback=teal!4,colframe=teal!35!black,title={Before you move on}]
What does \mw{मोकलणो} mean? (\textquotedblleft{}To send off\textquotedblright{}.) Say it once more.
\end{tcolorbox}
